\section{1D Residual Networks: Theoretical Background}
This section develops the theoretical foundation for 1D Residual Networks (ResNets) applied to time-series classification. We build directly on the convolutional machinery established in the 1D-CNN section, so the focus here is the residual learning principle, its motivation, and the normalization scheme used in our implementation.

\subsection{The Vanishing Gradient Problem in Deep CNNs}
Stacking convolutional blocks should, in principle, increase the network's expressive capacity. However, as the depth $L$ grows, the gradient propagated to the early layers shrinks geometrically. Recall from the 1D-CNN backpropagation derivation that the gradient of the loss with respect to the input of layer $\ell$ is obtained by chaining Jacobians:
\begin{equation}
    \frac{\partial \mathcal{L}}{\partial x_{\ell}} = \frac{\partial \mathcal{L}}{\partial x_{L}} \prod_{k=\ell}^{L-1} \frac{\partial x_{k+1}}{\partial x_{k}}
\end{equation}
Each Jacobian factor $\partial x_{k+1}/\partial x_{k}$ is a product of a convolution weight tensor and the derivative of the activation. With ReLU, this derivative is either 0 or 1, and the convolution weights are initialised with variance close to $1/\text{fan\_in}$. The product of many such terms collapses towards zero exponentially fast, freezing the early filters and preventing them from learning useful features. Beyond a depth of roughly 10 plain convolutional layers, training accuracy actually \textit{degrades} despite increased capacity --- a counter-intuitive phenomenon first documented by He et al.\ (2016).

\subsection{Residual Learning}
The residual block reformulates the problem. Instead of asking a stack of layers to learn a direct mapping $\mathcal{H}(x)$, we ask it to learn the \textbf{residual} $\mathcal{F}(x) = \mathcal{H}(x) - x$. The block's output is then constructed by an additive shortcut:
\begin{equation}
    y = \mathcal{F}(x; \mathbf{W}) + x
\end{equation}
The intuition is geometric. If the optimal mapping for a given block is the identity, a plain convolutional stack must painstakingly drive its weights to mimic identity through non-linearities. A residual block achieves the same result by simply driving $\mathcal{F}(x)$ to zero, which is the default behaviour at initialisation (weights centred on zero). This makes the identity an \textbf{attractor} of the optimisation landscape rather than a difficult target.

The crucial property is what the shortcut does to the backward pass. Differentiating the residual equation:
\begin{equation}
    \frac{\partial y}{\partial x} = \frac{\partial \mathcal{F}}{\partial x} + \mathbf{I}
\end{equation}
The identity term guarantees that the gradient flowing into the block is the sum of the gradient flowing through the convolutional body \textit{plus} an undamped copy that bypasses it entirely. Chained across $L$ blocks, the gradient at layer $\ell$ becomes:
\begin{equation}
    \frac{\partial \mathcal{L}}{\partial x_{\ell}} = \frac{\partial \mathcal{L}}{\partial x_{L}} \prod_{k=\ell}^{L-1} \left( \mathbf{I} + \frac{\partial \mathcal{F}_{k}}{\partial x_{k}} \right)
\end{equation}
Expanding the product yields a sum of terms, one of which is the bare $\partial \mathcal{L}/\partial x_{L}$ multiplied by 1. The early layers therefore always receive a usable error signal regardless of how deep the network grows, which is why ResNets can be trained to hundreds of layers without optimisation collapse.

\subsection{The Dimensional-Mismatch Shortcut}
The additive identity only works when $\mathcal{F}(x)$ and $x$ have the same number of channels. When a block increases the channel count (e.g.\ from 16 to 32), the shortcut must perform a learned dimensional projection. The standard solution is a $1 \times 1$ convolution:
\begin{equation}
    y = \mathcal{F}(x; \mathbf{W}) + \mathbf{W}_{s} \ast x
\end{equation}
where $\mathbf{W}_{s} \in \mathbb{R}^{C_{out} \times C_{in} \times 1}$ is a small learnable tensor with negligible parameter cost. The $1 \times 1$ convolution acts as a per-time-step linear projection of the channel dimension, preserving the temporal length so the shapes line up before summation.

\subsection{Group Normalization}
Our implementation replaces \texttt{BatchNorm1d} with \texttt{GroupNorm} in every block. BatchNorm normalises each channel using statistics aggregated across the batch dimension, then maintains a running mean and variance that are frozen at inference time. This creates a subtle but devastating problem in our setting: the validation and test batches contain a different mixture of anomaly classes than the training batches, so the frozen running statistics no longer match the actual feature distribution, producing erratic validation loss spikes that mask genuine convergence.

GroupNorm sidesteps the batch entirely. It partitions the $C$ channels of a feature map into $G$ groups (we use $G = \min(8, C)$) and normalises each group independently per-sample:
\begin{equation}
    \hat{x}_{c,t} = \frac{x_{c,t} - \mu_{g}}{\sqrt{\sigma_{g}^{2} + \varepsilon}}, \quad
    \mu_{g} = \frac{1}{|g| \cdot T} \sum_{c \in g} \sum_{t} x_{c,t}
\end{equation}
where $g$ is the group containing channel $c$. Because the statistics are computed from the sample itself, there are \textbf{no running statistics} and the train/eval behaviour is mathematically identical. The validation curves become monotone and trustworthy.

\newpage
